\documentclass[11pt]{amsart}
\usepackage[T1]{fontenc}
\usepackage[utf8]{inputenc}
\usepackage[french]{babel}
\usepackage{amssymb}
\usepackage[margin=2.5cm]{geometry}
\usepackage[hidelinks]{hyperref}
\newtheorem{proposition}{Proposition}
\newtheorem{corollary}[proposition]{Corollaire}
\title[Contrôle filtré : hypothèses et formalisation]{Contrôle fini d'un complexe filtré :\protect\\
hypothèses de génération et portée de la formalisation}
\date{27 octobre 2026}
\begin{document}
\maketitle

\noindent\textbf{Statut corrigé.}
Le développement \texttt{sag/filtered-cohomology} vérifie les implications
élémentaires à partir des hypothèses de génération (G) et (K) ci-dessous.
Il n'utilise \emph{aucune hypothèse de noethérianité ni de cohérence}.
Il ne construit pas (G) ou (K) pour les preuves ACL/ZMT ou étale-flat :
\emph{ce n'est donc pas encore un théorème commun instancié dans ces
deux preuves}. La première version de cette note ne distinguait pas assez
ces points. Les modules de cohomologie ne sont pas construits comme
quotients dans ce développement : on utilise des représentants et des
identités modulo un bord.

\section{Trois notions à distinguer}

\subsection*{Noethérianité classique}
Dans le rappel classique de la section suivante, « $A$ noethérien »
signifie que \emph{tout idéal de $A$ est de type fini}, dans le cadre
de la logique classique usuelle. On utilise aussi ses conséquences
classiques : un sous-module d'un module de type fini est de type fini,
et une algèbre de type fini sur $A$ est noethérienne.
Ces conséquences ne sont pas des hypothèses implicites du code nouveau.

\subsection*{Noethérianité par une barre dans Unit}
Dans \texttt{sag/noetherian/Span.seg}, une présentation comporte un type
de codes $B$ et une interprétation $\phi:B\to A$, compatible avec les
opérations de l'anneau. Une liste est \emph{bonne} si l'un de ses éléments
interprétés est une combinaison linéaire de ceux qui le précèdent,
avec coefficients eux-mêmes codés. Le plus récent élément est stocké
en tête de liste.

La donnée \texttt{Noeth} est exactement
\[
 \texttt{Bar.Over[B,Good].Bars(nil)}.
\]
Elle est constituée d'un arbre inductif bien fondé, ramifié sur $B$,
et d'une preuve que chaque feuille atteinte depuis la liste vide
porte une liste bonne. Une feuille est permise si la liste est bonne ;
un nœud est permis si l'on dispose d'un sous-arbre pour chaque
prolongement de la liste par un code. Il ne s'agit pas d'un arbre
nécessairement fini, ni d'une borne uniforme sur la longueur des branches.

Suivre une suite de codes dans cet arbre fournit une dépendance
\[
 \phi(b_N)=\sum_{i<N}\phi(c_i)\phi(b_i).
\]
La barre fournit également un principe d'induction : sa définition ne
se réduit pas à affirmer que chaque suite admet une telle dépendance.
\texttt{CodedNoetherian} ajoute que tout élément de l'anneau possède
un code sous troncation propositionnelle. La marche dans la barre
peut alors rester sous troncation, sans choisir une suite entière de codes.
C'est cette notion qui intervient dans les preuves existantes de
ZMT et de platitude.

\subsection*{Cohérence constructive}
Une formulation effective de la cohérence demande, pour chaque matrice
finie $D:A^r\to A^s$, une famille finie génératrice de $\ker D$, avec
la propriété de génération. Elle permet de calculer les syzygies et,
à partir de présentations finies, des présentations de noyaux.
La propriété de génération peut être exprimée sous troncation ;
elle ne demande pas de décider l'appartenance à un idéal arbitraire.
La caractérisation usuelle par les noyaux entre modules de présentation
finie est rappelée dans
\href{https://stacks.math.columbia.edu/tag/05CU}{Stacks, section 10.90}.

La cohérence est donc pertinente pour \emph{construire} certaines données
de génération. Mais elle ne dit pas qu'un sous-module arbitraire d'un
module de présentation finie est de type fini. Elle ne fournit pas
non plus la dépendance sur toute suite donnée qu'utilisent les preuves
existantes. La section~\ref{sec:coherence} précise l'obstacle ici.

\section{Le théorème classique et la distinction constructive}

Soient $A$ un anneau noethérien, $I$ un idéal, $C^\bullet$ un complexe,
et $q$ un degré. Posons
\[
 \mathcal R=\bigoplus_{n\geq0}I^nT^n,
 \qquad M_n=H^q(I^nC^\bullet),\qquad
 M=\bigoplus_{n\geq0}M_n.
\]
Si $M$ est de type fini sur $\mathcal R$, on obtient classiquement une
borne $c$ telle que, pour $n\geq c$,
\begin{align*}
 \operatorname{im}(M_n\longrightarrow M_0)&\subseteq I^{n-c}M_0,\\
 \ker(M_n\longrightarrow M_0)&\subseteq
 \ker(M_n\longrightarrow M_{n-c}).
\end{align*}
Le premier contrôle utilise des générateurs homogènes de $M$.
Le second utilise des générateurs homogènes du sous-module
\[
 K=\bigoplus_{n\geq0}\ker(M_n\longrightarrow M_0).
\]
Classiquement, $K$ est de type fini puisque $\mathcal R$ est noethérien
et $M$ de type fini. Ce contrôle se retrouve dans la preuve du
\href{https://stacks.math.columbia.edu/tag/02O7}{lemme 30.20.3 du Stacks Project}.

La finitude de $M$ est une \emph{hypothèse supplémentaire} ; elle ne
découle pas de la seule noethérianité de $A$ pour un complexe arbitraire.
Une fois $M$ et $K$ munis de familles génératrices homogènes, les deux
contrôles sont des calculs finis. Les sections suivantes formalisent
ces calculs, mais pas la construction de ces familles.

\section{Les classes représentables à un ordre élevé}

Écrivons $C^0\xrightarrow d C^1$ pour la différentielle qui précède
le degré de cohomologie étudié. Ainsi $C^1$ représente ici $C^q$ et
$C^0$ représente $C^{q-1}$. Une condition $Z(x)$ désigne les cycles
dans $C^1$.

Fixons une famille finie de cycles $z_i$, de degrés étiquetés $e_i$,
une borne $c$ et des entiers $g_i$ avec $e_i+g_i=c$.
\emph{Aucune hypothèse de noethérianité ou de cohérence n'est faite ici.}
La donnée de génération utilisée est la suivante : pour chaque $m\geq0$
et chaque cycle $x\in I^{c+m}C^1$, on dispose, éventuellement seulement
sous troncation, de $r_i$ et de $v$ tels que
\[
 r_i\in I^{g_i+m},\qquad v\in I^{c+m}C^0,
 \qquad x=\sum_i r_i z_i+dv.                 \tag{G}
\]
Si les $z_i$ représentent des générateurs homogènes de $M$, cette
donnée est exactement leur propriété de génération en degré $c+m$.
Dans cette interprétation, on doit aussi avoir $z_i\in I^{e_i}C^1$.
Le code ne requiert pas cette dernière propriété pour effectuer le calcul.

Le mot « certificat » désigne ici les coefficients, le bord et les
preuves des appartenances et de l'identité. Pour le théorème uniforme,
on suppose une \emph{preuve de l'existence d'un tel certificat pour
tout $m$ et tout $x$ concerné}, avec la même famille et la même borne $c$.
Il ne suffit donc pas de fournir une liste de candidats générateurs,
ni quelques exemples de décompositions. Sous troncation, on affirme
l'existence de la décomposition sans en choisir une globalement.
Cette propriété universelle est une hypothèse mathématique forte.

\begin{proposition}
Sous la donnée \textup{(G)}, la classe de $x$ dans $H^q(C^\bullet)$
appartient à $I^mH^q(C^\bullet)$.
\end{proposition}
\begin{proof}
Chaque $r_i$ appartient à $I^m$, puisque $g_i\geq0$.
L'identité (G) fournit donc le certificat
$[x]=\sum_i r_i[z_i]$ avec des coefficients dans $I^m$.
\end{proof}

Dans Unit : \texttt{FilteredControl.Control.Image.control} et
\texttt{UniformImage.image-control}. La conclusion formelle est
le certificat sur les représentants, sans quotient cohomologique.

\section{Les classes devenant nulles et le relèvement}

Fixons maintenant une famille finie de primitives $b_i\in C^0$,
toujours sans noethérianité ni cohérence.
Dans l'interprétation par le module $K$, on a $db_i\in I^{e_i}C^1$,
et $e_i$ est le degré de sa classe dans $H^q(I^{e_i}C^\bullet)$ ;
on ne suppose pas $b_i\in I^{e_i}C^0$.
De nouveau, $e_i+g_i=c$.

La donnée de génération du noyau est : pour tout $z\in C^0$ avec
$dz\in I^{c+m}C^1$, on dispose, éventuellement seulement sous
troncation, d'un certificat
\[
 r_i\in I^{g_i+m},\qquad v\in I^{c+m}C^0,
 \qquad dz=\sum_i r_i\,db_i+dv.              \tag{K}
\]
Des générateurs homogènes de $K$, avec représentants et primitives,
fournissent précisément une telle donnée.
Ici encore, la famille et $c$ sont fixes et la propriété est universelle
en $m,z$ : le code \texttt{UniformKernel} prend cette propriété comme
paramètre \texttt{generation}, il ne la démontre pas.

\begin{proposition}
Sous \textup{(K)}, il existe $w\in I^mC^0$ avec $dw=dz$.
Par conséquent il existe $h$ fermé avec $h-z\in I^mC^0$.
\end{proposition}
\begin{proof}
Poser $w=\sum_i r_i b_i+v$. Tous ses termes sont dans $I^mC^0$,
et (K) donne $dw=dz$. Prendre $h=z-w$.
\end{proof}

Dans Unit : \texttt{Kernel.control}, \texttt{UniformKernel.kernel-control}
et \texttt{UniformKernel.lift}. La génération sous troncation produit
un résultat sous troncation, sans choisir simultanément une suite de
représentants.

Cette formulation relève la section modulo $I^m$, dès que l'on
dispose d'une approximation à l'ordre $c+m$. Pour le lemme de correction
finie utilisé dans ACL, on prend $m=1$. L'énoncé uniforme demande
cependant plus de données que la seule dépendance de la suite des
obstructions effectivement rencontrées.

\section{L'annulation de la torsion infiniment divisible}

\begin{corollary}
Dans la situation de \textup{(G)}, supposons les bords nuls et $I$
annulant les cycles. Alors tout cycle d'ordre $c+m+1$ est nul.
\end{corollary}
\begin{proof}
Appliquer (G) avec $m+1$. Les coefficients $r_i$ appartiennent à $I$,
donc $r_i z_i=0$, et $dv=0$ par hypothèse.
\end{proof}

Dans Unit : \texttt{UniformImage.Kill.vanish}.

Pour décrire une application \emph{conditionnelle} liée au lemme Kill,
prendre un anneau $E$,
$I=(f)$, et le complexe $E\xrightarrow{f}E$ en degrés $0,1$.
On a
\[
 H^0(C^\bullet)=\operatorname{Ann}(f),\qquad
 H^0(I^nC^\bullet)=f^nE\cap\operatorname{Ann}(f).
\]
Dans la formalisation par représentants, on modélise les bords nuls
par l'application nulle $E\to E$, et les cycles par $fx=0$.
L'instance \texttt{TwoGenerated.Principal.Finite.vanish} donne
\[
 fx=0,\qquad x\in(f)^{c+m+1}\quad\Longrightarrow\quad x=0,
\]
sous les certificats de génération du module gradué d'annulation filtrée.
Les relèvements de Taylor de \texttt{etale-flat.tex} fournissent justement
un élément dans ces intersections pour tout ordre. Leur preuve actuelle
utilise une dépendance finie entre témoins et ne construit pas la borne
uniforme $c$ du présent corollaire.
En particulier, \texttt{TwoGenerated.Principal.Finite.vanish} conserve
le paramètre \texttt{generation} : il ne constitue pas une nouvelle
preuve du lemme Kill à partir de \texttt{CodedNoetherian}.

\section{La cohérence peut-elle fournir les hypothèses ?}\label{sec:coherence}

Oui, \emph{si l'on présente les modules concernés comme des noyaux de
morphismes entre modules de présentation finie sur un anneau cohérent}.
Avec des présentations effectives et une cohérence effective, on peut
alors obtenir des générateurs et leurs propriétés de génération.
Pour un noyau gradué, on peut ensuite prendre les composantes homogènes
des générateurs et majorer leurs degrés.

Mais ce travail de présentation n'est pas fait dans la note.
En particulier, écrire
\[
 K=\ker\left(M\longrightarrow H^q(C^\bullet)[T]\right)
\]
ne suffit pas : la cible est considérée comme module sur
$\mathcal R=\bigoplus I^nT^n$, et sa présentation finie sur cet anneau
n'est pas établie. La cohérence de $A$ ne donne pas à elle seule celle
de $\mathcal R$, ni les présentations nécessaires. De même, les termes
du complexe de \v Cech comportent des localisations ; ils ne sont pas
automatiquement des modules de présentation finie sur la base.

\subsection*{Pourquoi la cohérence seule ne suffit pas au relèvement}
Voici un contre-exemple algébrique, indépendant du code Unit.
Prenons dans $\mathbb Q(x,y)$ la valuation à valeurs dans
$\mathbb Z^2$ ordonné lexicographiquement, avec
\[
 \nu(x)=(0,1),\qquad \nu(y)=(1,0),
\]
et son anneau de valuation $V$. Concrètement, on prend d'abord l'ordre
en $y$, puis l'ordre en $x$ du coefficient initial en $y$.
L'anneau $V$ est cohérent : tout idéal de type fini non nul y est
principal et libre. Pour tout $n\geq0$, on a $y/x^n\in V$, donc
\[
 y\in x^nV\quad\hbox{pour tout }n.
\]
Dans le complexe $V\xrightarrow{\,y\,}V$, avec $I=(x)$ et $z=1$,
on a ainsi $dz\in I^nV$ à tout ordre. Pourtant, il n'existe aucun
$h$ avec $yh=0$ et $h\equiv1\pmod{xV}$ : l'intégrité force $h=0$,
alors que $x$ n'est pas inversible.

Ainsi, même pour un complexe de modules libres de rang un sur un
anneau cohérent, le relèvement à partir de tous les ordres peut échouer.
Ce contre-exemple réfute le remplacement de toute condition de finitude
filtrée par la seule cohérence de la base ; il ne réfute pas une approche
par cohérence assortie des présentations graduées appropriées.

\section{Applications : ce qui est déjà instancié et ce qui manque}

\subsection*{Pour ACL/ZMT}
La voie uniforme de cette note demanderait (K) pour le complexe de
\v Cech, avec une famille fixe de primitives et une borne valable pour
\emph{tous} les bords filtrés. Cette donnée n'est pas construite.

La preuve existante utilise moins. Elle représente la suite particulière
des obstructions dans une famille finie contrôlée en degré ;
\texttt{GradedDependence} fournit une relation
\[
 \rho_N=\sum_{1\leq i<N}r_i\rho_i,
 \qquad r_i\in I^{N-i}.
\]
\texttt{ReesLifting} et sa version modulo relations transforment cette
relation en une identité entre bords ; \texttt{FiniteCorrection}
en déduit la section corrigée. Cette voie ne demande pas des générateurs
de tout $K$. Les données de représentation de la suite sont effectivement
construites dans la preuve ACL, notamment dans \texttt{CechF1}.

\subsection*{Pour la platitude}
La voie uniforme demanderait (G) pour
$\bigoplus_n(f^nE\cap\operatorname{Ann}(f))$. Cette génération graduée
n'est pas construite non plus.

La preuve existante produit, pour l'élément $u$ étudié, des témoins
\[
 u=f^nw_n,\qquad f^{n+1}w_n=0\quad(n\geq0).
\]
La noethérianité codée d'un sous-anneau contenant ces données donne
$w_N=\sum_{i<N}\lambda_iw_i$. Alors
\[
 u=f^Nw_N=\sum_{i<N}\lambda_i f^{N-i-1}(f^{i+1}w_i)=0.
\]
C'est \texttt{CodedSubring.Vanish}, utilisé dans
\texttt{sag/etale-flat/RegularCoded.seg}. Le sous-anneau et les codes
sont construits dans cette preuve ; cela n'impose pas la noethérianité
à l'anneau de base arbitraire du théorème de platitude.

\subsection*{Réponse sur la portée de la nouvelle formalisation}
Les calculs conditionnels sont utilisables si l'on fournit (G) ou (K),
mais \emph{ils ne sont actuellement instanciés dans aucune de ces deux
preuves avec toutes leurs hypothèses déchargées}. Leur vérification
ne résout pas le problème de construire la génération filtrée.
Pour isoler un résultat commun aux preuves existantes, la voie déjà
justifiée est celle des dépendances sur les suites codées et de leur
utilisation dans une identité finie. Une autre voie, par cohérence,
devrait exhiber les présentations finies précises qui lui permettent
de construire (G) ou (K).

\section{Généralité et portée du code}

\texttt{ScalarFiltration} décrit des prédicats $J_n(r)$ décroissants.
\texttt{Filtration} décrit des sous-modules filtrés $F_n$ avec
$J_n(r)\Rightarrow rx\in F_n$, et les opérations finies nécessaires.
Le calcul abstrait s'applique donc aux puissances de tout idéal pour
lequel ces données sont fournies. \texttt{TwoGenerated.Over} fournit
l'instance déjà disponible dans la bibliothèque : $I=(a,b)$, y compris
les idéaux principaux en prenant $b=0$.

\texttt{Generator} contient la valeur et les étiquettes $e_i,g_i,c$,
avec une preuve de $e_i+g_i=c$. Ce record ne prétend pas certifier
à lui seul la présence de la valeur en degré $e_i$ dans un module de
Rees : le sens de la donnée est assuré par (G) ou (K). Pour la seconde
donnée, la valeur stockée est $b_i$, et l'étiquette concerne $db_i$.

Le résultat vérifié est donc un théorème uniforme \emph{à partir de
certificats de génération}. Il reste à construire ces certificats
pour les complexes concernés si l'on souhaite remplacer les preuves
existantes par cet énoncé, ou à formaliser des hypothèses effectives
suffisantes pour construire les générateurs de $M$ et de $K$.
Les preuves ACL et étale-flat existantes n'ont pas été modifiées.

\end{document}
